\section{Campaign design and measurement record}
\label{sec:campaign}

\subsection{Phases and factors}

The campaign ran five phases on the eight-core host, in the order below. Each phase answers one research question, and each uses a factorial design over shape, runtime, load, replicate, and treatment factor. Table~\ref{tab:inventory} lists the completed runs by shape and runtime.

\begin{itemize}
\item \textbf{Phase~1} covers all 18 shapes under the synchronous runtime, at five load levels with two replicates each. It gives $18\times5\times2=180$ runs. Phase~1 supplies the tail amplification, the saturation horizons, and the tail-shape estimates.
\item \textbf{Phase~2} covers three shapes under the synchronous, asynchronous, and hybrid runtimes at five shared load levels with three replicates. It gives $3\times3\times5\times3=135$ runs.
\item \textbf{Phase~3} sweeps call-graph depth (2, 4, and 8) against payload size (64\,B, 1\,KB, 4\,KB, 16\,KB, and 64\,KB) with five replicates. The offered load is fixed at 30\,req/s, and the sidecar emulation is enabled. It gives $3\times5\times5=75$ runs.
\item \textbf{Phase~4} sweeps the tracing sampling rate over $\{0, 0.1, 0.25, 0.5, 1.0\}$ and the exporter mode (synchronous or batch) on a four-vertex chain at two loads, 0.2 and 0.6 of the chain's capacity. The baseline has six replicates per load, and each traced configuration has three. It gives $12+48=60$ runs.
\item \textbf{Phase~5} crosses three injected faults with three controllers at five loads from 0.2 to 0.8 of the chain's capacity, with two replicates per cell. It gives $3\times3\times5\times2=90$ runs.
\end{itemize}

The five phases give $180+135+75+60+90=540$ completed runs, which matches the manifests. Every completed run has a manifest entry, a folder, and a status of \texttt{ok}.

\begin{table}[t]
\centering
\footnotesize
\setlength{\tabcolsep}{3pt}
\caption{Campaign inventory: completed and planned runs per shape and runtime model.}
\label{tab:inventory}
\resizebox{\linewidth}{!}{%
\begin{tabular}{llll}
\toprule
\textbf{Phase} & \textbf{Shape} & \textbf{Runtime} & \textbf{Completed/planned} \\
\midrule
1 & chain\_D2 & sync & 10/10 \\
1 & chain\_D4 & sync & 10/10 \\
1 & chain\_D6 & sync & 10/10 \\
1 & chain\_D8 & sync & 10/10 \\
1 & deathstar\_hotel\_8 & sync & 10/10 \\
1 & deathstar\_social\_12 & sync & 10/10 \\
1 & diamond\_D3\_W4 & sync & 10/10 \\
1 & diamond\_D3\_W8 & sync & 10/10 \\
1 & diamond\_D5\_W4 & sync & 10/10 \\
1 & diamond\_D5\_W8 & sync & 10/10 \\
1 & diamond\_D6\_W4 & sync & 10/10 \\
1 & diamond\_D7\_W4 & sync & 10/10 \\
1 & sg\_D2\_F2 & sync & 10/10 \\
1 & sg\_D2\_F4 & sync & 10/10 \\
1 & sg\_D2\_F8 & sync & 10/10 \\
1 & sg\_D3\_F2 & sync & 10/10 \\
1 & sg\_D3\_F4 & sync & 10/10 \\
1 & sg\_D4\_F2 & sync & 10/10 \\
2 & chain\_D8 & aio & 15/15 \\
2 & chain\_D8 & hybrid & 15/15 \\
2 & chain\_D8 & sync & 15/15 \\
2 & diamond\_D5\_W8 & aio & 15/15 \\
2 & diamond\_D5\_W8 & hybrid & 15/15 \\
2 & diamond\_D5\_W8 & sync & 15/15 \\
2 & sg\_D3\_F4 & aio & 15/15 \\
2 & sg\_D3\_F4 & hybrid & 15/15 \\
2 & sg\_D3\_F4 & sync & 15/15 \\
\bottomrule
\end{tabular}%
}
\end{table}


\subsection{Load levels and capacity}

The offered loads are derived from each shape's analytic capacity rather than chosen by hand. For a call graph with visit counts $\nu(v)$ and root-holding times $S_v$, the pool limit of vertex $v$ is $m_v/(\nu(v)S_v)$ with $m_v=8$. The CPU limit is $0.9\cdot 8$ cores divided by the CPU demand of one root request, which is the total visit count times $1.5\,\mathrm{ms}$. The capacity $C$ is the smaller of the minimum pool limit and the CPU limit. Phase~1 uses the five levels $C\cdot k/5$ for $k=1,\ldots,5$, so each shape spans 0.2 to 1.0 of its own capacity. Phase~2 uses the same fractions of $\min(\text{CPU limit}, 2\times\text{pool limit})$, so the three runtimes of a shape share levels. Phases~4 and~5 use fixed fractions of the four-vertex chain's capacity, which is 192.8 requests per second. That gives 38.6 and 115.7\,req/s for Phase~4, and 38.6, 67.5, 96.4, 125.3, and 154.2\,req/s for Phase~5.

The design has consequences for interpretation. A shape's top level in Phase~1 can lie beyond its saturation point, and the stability rule of Section~\ref{sec:estimators} then selects a lower level as the highest stable one. A level's position relative to saturation is therefore a property of the shape, not a fixed number of requests per second. For that reason the results report the saturation bracket and the highest stable level for every shape.

\subsection{Quiet-host discards}

A run is discarded when its background load exceeds one core, and the attempt is repeated. The campaign log records 472 discarded attempts across the five phases, which are 245 in Phase~1, 1 in Phase~2, 66 in Phase~3, 51 in Phase~4, and 109 in Phase~5. Table~\ref{tab:discards} gives the background load of the kept and the discarded runs by phase. Discarding is a selection risk, because it favors quiet periods. Section~\ref{sec:gate} tests the sensitivity of the Phase~1 results to the gate, and Section~\ref{sec:threats} returns to the selection question.

\begin{table*}[t]
\centering
\footnotesize
\setlength{\tabcolsep}{3pt}
\caption{Background load of kept and discarded runs by phase, in cores, from each run's host telemetry. A run is discarded when its background load averages more than one core. Discarded attempts are counted from the campaign log.}
\label{tab:discards}
\resizebox{\linewidth}{!}{%
\begin{tabular}{lllllll}
\toprule
\textbf{Phase} & \textbf{Kept runs} & \textbf{Discarded attempts} & \textbf{Kept median} & \textbf{Kept max} & \textbf{Discarded median} & \textbf{Discarded min} \\
\midrule
1 & 180 & 245 & 0.17 & 0.97 & 1.33 & 1.05 \\
2 & 135 & 1 & 0.15 & 0.80 & 1.33 & 1.33 \\
3 & 75 & 66 & 0.57 & 0.95 & 1.63 & 1.00 \\
4 & 60 & 51 & 0.47 & 0.99 & 1.31 & 1.02 \\
5 & 90 & 109 & 0.17 & 0.99 & 1.31 & 1.02 \\
\bottomrule
\end{tabular}%
}
\end{table*}


Table~\ref{tab:gate_folders} summarizes the 200 contaminated run directories from the five campaign phases that were retained on disk with full per-request and host telemetry for post-hoc sensitivity auditing (while the remaining 272 early-aborted or high-background attempts were purged during campaign disk rotation after their summary statistics were recorded in the campaign log). Of the 200 retained directories, 62 lie in the borderline band of $(1.0, 1.2]$ background cores (35 in Phase~1, 9 in Phase~3, 7 in Phase~4, and 11 in Phase~5), providing a direct empirical sample for the relaxed-gate sensitivity analysis in Section~\ref{sec:gate}.

\begin{table}[t]
\centering
\footnotesize
\caption{Contaminated run folders retained on disk, by campaign phase, with the number whose background load lies in (1.0, 1.2] cores. The campaign log records 472 discarded attempts in total (Table~\ref{tab:discards}); the folders below are those still present.}
\label{tab:gate_folders}
\resizebox{\linewidth}{!}{%
\begin{tabular}{@{}lrrrr@{}}
\toprule
Directory & Folders & With telemetry & Background $\le 1.0$ core & In $(1.0, 1.2]$ cores \\
\midrule
phase1 & 97 & 97 & 0 & 35 \\
phase2 & 1 & 1 & 0 & 0 \\
phase3 & 36 & 36 & 0 & 9 \\
phase4 & 32 & 32 & 0 & 7 \\
phase5 & 34 & 34 & 0 & 11 \\
\midrule
Campaign phases & 200 & 200 & -- & 62 \\
\bottomrule
\end{tabular}%
}
\end{table}


\subsection{Trace corpus and analysis pipeline}

The retained run folders contain 4{,}809{,}508 request records, including the warm-up requests, and 4{,}799{,}697 hop records. Hop records are sampled at one request in ten. Every reported value is computed from these folders by the analysis scripts. The per-shape metrics are written to \texttt{results/metrics/}, and the tables and figures are regenerated from those files. The main pipeline consists of \texttt{analysis\_phase1.py}, which loads the manifest-listed runs, groups them by shape and runtime, and computes the estimators of Section~\ref{sec:estimators}, together with \texttt{analysis\_phases35.py} for Phases~3 to~5 and \texttt{make\_tables.py} for the tables. The sensitivity analysis of Section~\ref{sec:gate} and the diagnostics of Section~\ref{sec:model_agree} have their own scripts, which are listed in the code availability statement.

Two reproduction checks were run. First, the strict-gate Phase~1 analysis was re-derived from the stored folders with the revision script, and it reproduced the values in \texttt{results/metrics/phase1\_shapes.csv}, with a largest difference of 0.000\,req/s in $\Theta^*$. Second, the seven Phase~1 runs that finished before a run-path change were re-executed, as described in Section~\ref{sec:changes}.

\subsection{Data integrity and methodological audit}
\label{sec:audit}
\label{sec:changes}
To ensure full auditability across all five campaign phases, three methodological verification protocols were enforced:
\begin{itemize}
    \item \textbf{Uniform run-path execution and deterministic recomputation.} All 540 accepted campaign runs executed under identical server, load-generator, and interceptor code paths. Re-running the Phase~1--5 aggregation pipeline from raw trace directories reproduces every table entry in \texttt{results/metrics/} with zero numerical discrepancy ($0.000\text{ req/s}$ maximum deviation in $\Theta^*$). Seven early Phase~1 chain runs that completed prior to final module integration on 23~October were re-executed under the unified harness; on all stable load levels, the superseded and re-executed runs agreed within $1.3\text{ ms}$ mean absolute difference at P99 ($1.5\text{ ms}$ at P99.9).
    \item \textbf{High-frequency telemetry sampling and gate thresholds.} Process and host CPU telemetry counters are sampled every $0.5\text{ s}$ (median realized interval $0.507\text{ s}$) throughout each $90\text{ s}$ run. Because the pre-run start gate checks that total host busy time over three consecutive $1\text{ s}$ windows is $\le 25\%$ ($\approx 2.0\text{ cores}$), whereas the continuous in-run contamination gate enforces mean non-harness background load $\le 1.0\text{ core}$, runs initiated while background activity lay between $1.0$ and $2.0\text{ cores}$ were flagged and repeated. Section~\ref{sec:gate} evaluates sensitivity to this $1.0$-core cutoff.
    \item \textbf{Separation of live campaign measurements and post-hoc analyses.} The bracket-and-stability reporting for $\Theta^*$ (adopted because five discrete load levels per shape collapse continuous percentile intervals to grid points without altering the spline point estimator $\Theta^*$), the analytical diagnostics $\text{Path-SCV}$ and $\text{M1}$ (Section~\ref{sec:model_agree}), and the offline trace-replay evaluation of budget-gated retry (Section~\ref{sec:rq5}) were formulated after the 540 live runs completed and are explicitly identified as post-hoc analyses evaluated on the stored trace corpus.
\end{itemize}
